\chapter{White Blood Cell Differential}

\section*{What Is This Test?}
The white blood cell (WBC) differential, also known as the leukocyte differential count, measures the percentage and absolute number of each type of white blood cell in the peripheral blood. The five major leukocyte types quantified are neutrophils (including bands, the less mature form), lymphocytes, monocytes, eosinophils, and basophils. The differential may also identify and enumerate abnormal or immature cells, such as blast cells, promyelocytes, myelocytes, metamyelocytes, atypical lymphocytes, and nucleated red blood cells. The differential is a core component of the complete blood count (CBC with differential) and provides essential information for the diagnosis and monitoring of infections (bacterial, viral, parasitic), inflammatory conditions, allergic disorders, hematologic malignancies (leukemias, lymphomas), bone marrow disorders, and the hematologic effects of medications and toxins. Results are reported both as a relative percentage and as an absolute count (cells per microliter), with the absolute count being clinically more important for interpretation. The differential can be performed by automated hematology analyzers using flow cytometry with laser light scatter and fluorescence, or manually by microscopic examination of a stained peripheral blood smear.

\section*{Alternative Names}
WBC Differential; Leukocyte Differential Count; White Blood Cell Differential Count; CBC with Differential; Diff; Manual Differential; Automated Differential; Five-Part Differential

\section*{Principle}
On modern automated hematology analyzers (Sysmex XN, Beckman Coulter DxH, Siemens ADVIA), the WBC differential is determined using flow cytometric techniques. Red blood cells are lysed, and the remaining white blood cells are differentially labeled with fluorescent dyes or characterized by their physical properties. In Sysmex analyzers, a fluorescent polymethine dye binds to nucleic acids in white blood cells, and laser light scatter (forward scatter for cell size, side scatter for internal complexity) combined with fluorescence intensity (correlating with nucleic acid content) allows precise classification into the five major types plus immature granulocytes \cite{briggs2012performance}. Beckman Coulter uses VCS technology (volume, conductivity, scatter), where impedance measures cell volume, radiofrequency conductivity assesses internal complexity, and laser scatter evaluates surface structure. Siemens ADVIA employs peroxidase cytochemistry (neutrophils, eosinophils, and monocytes have peroxidase activity and are differentiated by peroxidase staining intensity and cell size) combined with laser light scatter for basophils (resistant to a different lysing reagent). Automated analyzers count 10,000--30,000 white blood cells for the differential, providing superior precision compared to the manual 100--200 cell count. However, automated analyzers cannot identify all morphologic abnormalities (Auer rods, toxic granulation, D\"{o}hle bodies, blast cells not fitting automated algorithms, plasma cells), so samples flagged by the analyzer must be reviewed manually. In manual differential counting, a Wright-Giemsa-stained peripheral blood smear is examined microscopically under oil immersion (100$\times$ objective), and 100--200 white blood cells are classified into their respective types based on nuclear and cytoplasmic morphology.

\section*{History}
The ability to differentiate white blood cell types began with Paul Ehrlich in the 1870s--1880s. Ehrlich developed aniline dye staining techniques that distinguished neutrophils (granules that stain with neutral dyes), eosinophils (granules with affinity for acid dyes like eosin), and basophils (granules with affinity for basic dyes). He also described mononuclear cells (lymphocytes and monocytes). The Romanowsky stain (1891) and its modifications by Wright (1902) and Giemsa (1904) became the standard staining methods, producing the characteristic differential coloration that enables manual differential counting. For most of the 20th century, the manual differential, performed by a medical technologist counting 100 cells under oil immersion microscopy, was a standard hematology procedure in every clinical laboratory. The development of automated five-part differential analyzers began in the 1970s with the Technicon Hemalog D (using cytochemistry and light scatter) and advanced rapidly through the 1980s and 1990s. Modern automated analyzers can now perform five-part differentials plus immature granulocyte and nucleated red blood cell counts in less than a minute, with precision far exceeding manual methods. Digital morphology systems (CellaVision, Medica) that scan and pre-classify cells for technologist review represent the most recent advancement.

\section*{Why This Test Is Ordered}
\begin{itemize}
  \item Diagnosis of bacterial infection --- neutrophilia with left shift (increased bands and immature myeloid cells), toxic granulation, and D\"{o}hle bodies suggest acute bacterial infection
  \item Diagnosis of viral infection --- lymphocytosis, often with atypical (reactive) lymphocytes, is characteristic of viral infections such as Epstein-Barr virus/infectious mononucleosis, cytomegalovirus, HIV, and viral hepatitis
  \item Diagnosis of parasitic and allergic conditions --- eosinophilia (absolute eosinophil count >500/$\mu$L) is associated with parasitic infections (especially helminths such as Strongyloides, Ascaris, hookworm), allergic disorders (asthma, atopic dermatitis, allergic rhinitis), and drug hypersensitivity reactions
  \item Diagnosis and classification of leukemia --- presence of blast cells in the peripheral blood requires immediate investigation; specific features such as Auer rods (acute myeloid leukemia), smudge cells (chronic lymphocytic leukemia), and hairy cells (hairy cell leukemia) are diagnostic
  \item Monitoring chemotherapy and bone marrow recovery --- serial differentials track resolution of leukopenia, appearance of immature myeloid cells indicating recovery, and return to normal differential patterns
  \item Evaluation of unexplained fever, weight loss, night sweats, fatigue, or lymphadenopathy
  \item Monitoring immunosuppressive therapy --- absolute neutrophil count (ANC) is used to assess risk of infection and guide treatment decisions in patients on chemotherapy or immunosuppressive drugs. ANC <500/$\mu$L indicates severe neutropenia with high infection risk
  \item Screening for congenital immunodeficiency --- absent or markedly low lymphocytes may indicate severe combined immunodeficiency (SCID); absent neutrophils suggests congenital agranulocytosis (Kostmann syndrome)
  \item Diagnosis and monitoring of myelodysplastic syndromes --- persistent cytopenias with dysplastic features on differential
  \item Preoperative assessment --- part of comprehensive preoperative evaluation
\end{itemize}

\section*{Primary vs Secondary Test}
The WBC differential is a primary component of the CBC with differential and provides essential information for the diagnosis and monitoring of infectious, inflammatory, allergic, and hematologic disorders.

\section*{How This Test Is Performed}
A venous blood sample is collected in a lavender-top (EDTA) tube, typically 2--5 mL from the antecubital fossa. The tube is gently inverted 8--10 times to ensure adequate mixing with the EDTA anticoagulant. For automated differential, the sample is placed on an automated hematology analyzer, which lyses the red blood cells, differentially stains or characterizes the white blood cells, and classifies them using flow cytometric methods. Results are generated within 1--2 minutes. For a manual differential, a peripheral blood smear is prepared on a glass slide, fixed in methanol, stained with Wright-Giemsa stain, and examined under oil immersion microscopy (100$\times$ objective). The technologist systematically scans the smear and identifies 100--200 white blood cells by their nuclear and cytoplasmic features, reporting the percentage of each cell type. Any abnormal cells, morphological abnormalities, or red blood cell/platelet morphologic findings are also reported. Manual differential typically requires 15--45 minutes of technical time. Automated differentials flagged for abnormal or immature cell populations are reflexively reviewed with a manual differential.

\section*{What Preparation Is Needed}
No special preparation is required. Fasting is not necessary for a WBC differential. However, patients should avoid vigorous exercise for 30 minutes before testing, as exercise can transiently increase the neutrophil count (demargination from vessel walls) and lymphocyte count. Medications that can affect the differential, particularly corticosteroids (increase neutrophils, decrease lymphocytes and eosinophils), lithium (increases neutrophils), epinephrine (demarginates neutrophils), antipsychotics (can cause agranulocytosis), and chemotherapy, should be disclosed. Recent infection, vaccination, or surgery can transiently alter the differential.

\section*{Contraindications and Precautions}
No absolute contraindications. Important pre-analytical and analytical considerations include: EDTA is the standard anticoagulant; the sample must be thoroughly mixed to prevent microclots which can trap white blood cells and falsely reduce counts; smears for manual differential should be prepared within 2--4 hours of collection to prevent degenerative morphological changes; hemolyzed or clotted specimens are unacceptable; prolonged EDTA storage causes neutrophilic nuclear degradation, vacuolization (can mimic toxic change), and lymphocyte nuclear distortion; marked thrombocytosis or leukocytosis can produce spurious automated differential results; cold agglutinins may cause white cell clumping; very low white blood cell counts (leukopenia) reduce accuracy of both automated and manual differentials (buffy coat smear preparation may be needed for manual review); nucleated red blood cells interfere with automated WBC count and may require manual correction. Automated analyzers produce flags (alarms) when abnormal cell populations are detected --- these flagged samples require manual smear review before results are finalized.

\section*{Risks}
Risks are those of routine venipuncture: mild pain at the needle site, bruising or hematoma, lightheadedness or vasovagal syncope, and very rarely, infection or nerve injury. No additional risks are associated specifically with the differential test beyond those of venipuncture.

\section*{What Happens After the Test}
Automated differential results are available within 30--60 minutes. Manual differential review may take an additional 1--4 hours if the automated result is flagged for abnormality. The healthcare provider reviews the absolute counts of each cell type, not just the percentages. Critical values that may trigger immediate notification include: absolute neutrophil count <500/$\mu$L (severe neutropenia, high infection risk), presence of blast cells (suspected acute leukemia), and absolute neutrophil count >100,000/$\mu$L in a new patient (possible leukemia). The differential result is interpreted in the context of the complete CBC, clinical history, physical examination, and any available peripheral smear morphology report. Abnormal results guide further testing --- flow cytometry for suspected leukemia, bone marrow biopsy for unexplained cytopenias or blasts, infectious disease serologies for atypical lymphocytosis, and allergy evaluation for eosinophilia.

\section*{Understanding Results}

\subsection*{Below Normal (Decreased Cell Types)}
\begin{itemize}
  \item \textbf{Neutropenia (ANC <1,500/$\mu$L):} Caused by viral infections (influenza, hepatitis, HIV, EBV), sepsis, drug-induced (chemotherapy, clozapine, sulfonamides, antithyroid medications, phenytoin), autoimmune neutropenia, congenital neutropenia (Kostmann syndrome, cyclic neutropenia), aplastic anemia, myelodysplastic syndromes, vitamin B12 or folate deficiency, hypersplenism, and bone marrow infiltration. Severe neutropenia (ANC <500/$\mu$L) carries high risk of life-threatening bacterial and fungal infections.
  \item \textbf{Lymphopenia (absolute lymphocytes <1,000/$\mu$L in adults):} Associated with HIV/AIDS (CD4+ T-cell depletion), corticosteroid use (redistribution of lymphocytes), severe acute infections (COVID-19, sepsis, miliary tuberculosis), systemic lupus erythematosus, radiation therapy, chemotherapy, aplastic anemia, Hodgkin lymphoma (advanced), and congenital immunodeficiency (SCID, DiGeorge syndrome, Wiskott-Aldrich syndrome).
  \item \textbf{Monocytopenia (absolute monocytes <200/$\mu$L):} Seen in hairy cell leukemia, aplastic anemia, corticosteroid therapy, and severe sepsis. Monocytopenia is a characteristic finding in hairy cell leukemia and often key to its diagnosis.
  \item \textbf{Eosinopenia (<50/$\mu$L):} Caused by acute bacterial infection, corticosteroid administration (suppresses eosinophil release from bone marrow), Cushing syndrome, surgery, trauma, and burns. Eosinopenia is a nonspecific indicator of acute stress or inflammation.
  \item \textbf{Basopenia (<20/$\mu$L):} Difficult to detect due to normally low basophil counts. Associated with corticosteroid use, hyperthyroidism, acute infection, and allergic reactions.
\end{itemize}

\subsection*{Above Normal (Increased Cell Types)}
\begin{itemize}
  \item \textbf{Neutrophilia (absolute neutrophils >7,000/$\mu$L):} Most commonly caused by acute bacterial infection (pneumonia, abscess, pyelonephritis, meningitis, sepsis), acute inflammation (burns, trauma, surgery, gout, myocardial infarction), tissue necrosis, corticosteroid use (demargination of neutrophils from vessel walls), lithium therapy, epinephrine administration, pregnancy (physiological), heavy exercise, smoking, and hematologic malignancies --- chronic myeloid leukemia (CML, markedly elevated neutrophils with left shift, basophilia, low leukocyte alkaline phosphatase score), polycythemia vera, and essential thrombocythemia.
  \item \textbf{Lymphocytosis (absolute lymphocytes >4,000/$\mu$L in adults):} Caused by viral infections --- Epstein-Barr virus/infectious mononucleosis (atypical lymphocytes, positive heterophile antibody/Monospot test), cytomegalovirus, adenovirus, hepatitis viruses, HIV (acute seroconversion), pertussis (Bordetella pertussis, marked lymphocytosis in children), brucellosis, toxoplasmosis, and hematologic malignancies --- chronic lymphocytic leukemia (CLL, monoclonal B-cell lymphocytosis with smudge cells, diagnosed by flow cytometry demonstrating CD5+/CD19+/CD23+ B cells), acute lymphoblastic leukemia (ALL, blast cells, most common childhood leukemia), non-Hodgkin lymphomas with leukemic phase (mantle cell lymphoma, follicular lymphoma, splenic marginal zone lymphoma), and large granular lymphocytic leukemia.
  \item \textbf{Monocytosis (absolute monocytes >800/$\mu$L):} Associated with recovery phase of acute bacterial infection, chronic infections --- tuberculosis, subacute bacterial endocarditis, brucellosis, syphilis, rickettsial infections, fungal infections, protozoal infections (malaria, leishmaniasis), chronic inflammatory conditions --- sarcoidosis, Crohn's disease, ulcerative colitis, systemic lupus erythematosus, rheumatoid arthritis, and hematologic malignancies --- chronic myelomonocytic leukemia (CMML, persistent monocytosis >1,000/$\mu$L for >3 months with dysplasia), acute monocytic and myelomonocytic leukemias (AML subtypes M4 and M5), Hodgkin lymphoma, and certain non-Hodgkin lymphomas.
  \item \textbf{Eosinophilia (absolute eosinophils >500/$\mu$L):} Classified as mild (500--1,500/$\mu$L), moderate (1,500--5,000/$\mu$L), or severe (>5,000/$\mu$L --- hypereosinophilic syndrome). Causes include: parasitic infections (helminths --- Strongyloides stercoralis, Ascaris lumbricoides, hookworm, Trichinella spiralis, Toxocara canis, Schistosoma species, filariasis), allergic disorders (asthma, atopic dermatitis, allergic rhinitis, urticaria, food allergy), drug hypersensitivity (antibiotics --- penicillins, cephalosporins, sulfonamides; NSAIDs; anticonvulsants --- phenytoin, carbamazepine; allopurinol), dermatologic conditions (eczema, pemphigus, dermatitis herpetiformis), pulmonary disorders (allergic bronchopulmonary aspergillosis, eosinophilic pneumonia, L\"{o}ffler syndrome, Churg-Strauss syndrome/eosinophilic granulomatosis with polyangiitis), gastrointestinal disorders (eosinophilic esophagitis, gastroenteritis, inflammatory bowel disease), adrenal insufficiency (Addison disease --- loss of glucocorticoid suppression of eosinophils), hematologic malignancies (chronic eosinophilic leukemia, acute myeloid leukemia with eosinophilia --- AML-M4eo with inv(16), Hodgkin lymphoma, T-cell lymphomas, systemic mastocytosis), and hypereosinophilic syndrome (persistent eosinophilia >1,500/$\mu$L with end-organ damage to heart, lungs, nervous system).
  \item \textbf{Basophilia (absolute basophils >100/$\mu$L):} Most commonly associated with myeloproliferative neoplasms --- chronic myeloid leukemia (basophilia is a hallmark, helping distinguish CML from reactive leukocytosis), polycythemia vera, essential thrombocythemia, primary myelofibrosis, and systemic mastocytosis. Also seen in hypothyroidism (myxedema), chronic hypersensitivity states, ulcerative colitis, and iron deficiency anemia (rare). Persistent unexplained basophilia should prompt investigation for myeloproliferative neoplasm.
\end{itemize}

\section*{Normal Results}
\begin{itemize}
  \item \textbf{Neutrophils (segmented + bands):} Relative 40--70\%; Absolute 1,800--7,800/$\mu$L (1.8--7.8 $\times$ 10\textsuperscript{9}/L)
  \item \textbf{Lymphocytes:} Relative 20--40\%; Absolute 1,000--4,800/$\mu$L (1.0--4.8 $\times$ 10\textsuperscript{9}/L)
  \item \textbf{Monocytes:} Relative 2--8\%; Absolute 200--800/$\mu$L (0.2--0.8 $\times$ 10\textsuperscript{9}/L)
  \item \textbf{Eosinophils:} Relative 1--4\%; Absolute 0--500/$\mu$L (0--0.5 $\times$ 10\textsuperscript{9}/L)
  \item \textbf{Basophils:} Relative 0--1\%; Absolute 0--150/$\mu$L (0--0.15 $\times$ 10\textsuperscript{9}/L)
  \item \textbf{Bands (immature neutrophils):} 0--6\% of total WBC; Absolute 0--700/$\mu$L
  \item \textbf{Newborns (first week):} Neutrophils 50--70\%, Lymphocytes 25--35\% --- neutrophil predominance at birth
  \item \textbf{Children (1--6 years):} Neutrophils 25--45\%, Lymphocytes 45--65\% --- lymphocyte predominance is normal
  \item \textbf{Children (6--12 years):} Transition to adult pattern; Neutrophils 35--55\%, Lymphocytes 35--50\%
  \item \textbf{No blast cells, promyelocytes, myelocytes, metamyelocytes, atypical lymphocytes, or nucleated red blood cells} should be present in normal peripheral blood
\end{itemize}

\section*{Follow-up Tests}
\begin{itemize}
  \item \textbf{Peripheral blood smear (manual review):} Essential when automated differential is flagged for abnormalities, blast cells, or atypical cells
  \item \textbf{Flow cytometry/immunophenotyping:} For suspected leukemia or lymphoma --- determines cell lineage (myeloid vs. lymphoid, B-cell vs. T-cell), clonality, and specific diagnostic markers (CD5/CD19 in CLL, CD10 in ALL, CD34 in AML)
  \item \textbf{Bone marrow aspiration and biopsy:} When blast cells, unexplained cytopenias, or suspected myelodysplasia, leukemia, or marrow infiltration are present
  \item \textbf{Cytogenetic analysis (karyotype) and FISH:} For suspected hematologic malignancy --- Philadelphia chromosome t(9;22) for CML, t(15;17) for acute promyelocytic leukemia, del(5q) for MDS
  \item \textbf{Molecular genetic testing:} BCR-ABL1 fusion gene (CML), JAK2 V617F mutation (myeloproliferative neoplasms), PML-RARA (APL)
  \item \textbf{Heterophile antibody test (Monospot) and EBV serologies:} For atypical lymphocytes and suspected infectious mononucleosis
  \item \textbf{HIV antibody/antigen testing:} For unexplained lymphopenia or lymphocytosis
  \item \textbf{Blood cultures and infectious disease serologies:} For suspected bacterial, viral, parasitic, or rickettsial infection
  \item \textbf{Stool ova and parasite examination and serologies for Strongyloides, Toxocara, Trichinella:} For unexplained eosinophilia
  \item \textbf{Serum immunoglobulins (IgG, IgA, IgM, IgE) and serum protein electrophoresis:} For suspected immunodeficiency, hypergammaglobulinemia, or plasma cell dyscrasia
  \item \textbf{CRP, ESR, ANA, rheumatoid factor:} For evaluation of underlying inflammatory or autoimmune disease
  \item \textbf{Allergy testing (skin prick testing, allergen-specific IgE):} For persistent eosinophilia with allergic symptoms
  \item \textbf{Leukocyte alkaline phosphatase (LAP) score:} To differentiate CML (low LAP score) from leukemoid reaction (high LAP score)
\end{itemize}

\section*{References}
\bibliographystyle{unsrtnat}
\bibliography{references}

\clearpage
\section*{Regulatory Status and Authorization Date}
Regulatory status applies to the specific assay, instrument, manufacturer, and intended use rather than to the test name alone. No single approval date can be assigned to this test category. For a commercial assay, verify the current product in the relevant regulator's database: the U.S. Food and Drug Administration (FDA) clearance, approval, or authorization; India's Central Drugs Standard Control Organisation (CDSCO) licence or permission; the European Union In Vitro Diagnostic Regulation (EU IVDR) CE certificate issued through the applicable conformity-assessment route; or the United Kingdom Medicines and Healthcare products Regulatory Agency (MHRA) registration and UKCA/CE status. Laboratory-developed tests may not have an individual FDA, CDSCO, EU IVDR, or MHRA approval date.
\clearpage
